\documentclass[12pt,a4paper]{article}

\usepackage[hidelayout]{../../../Latex/Style/coursework-report-core}

\addbibresource{res-dat.bib}

\title{RES-DAT --- Environmental Data and Case-Study Definition}
\author{}
\date{\today}

\begin{document}

\maketitle

\tableofcontents
\clearpage

\section{Domain Purpose and Scope}
\label{sec:res-dat-domain-purpose}

RES-DAT defines the environmental data and case-study evidence needed to
reconstruct one selected historical tsunami event, validate that event against
authoritative observations, and support documented defence-specific impact
cases. The initial project retains one primary historical event and controlled
physical variations around that event. It does not add unrelated tsunami events
or build a broad multi-event ML dataset during the initial scope.

The active studentship case is the 2011 Tōhoku event with two
regional-to-local corridor trajectories: epicentre to the Kamaishi
region and epicentre to the Sendai coastal plain. The active use of the
data is impact modelling, Tōhoku/local comparison and barrier-class
comparison. ML datasets, surrogate optimisation and generative geometry
remain future work until the event reconstruction, local hydrodynamic
load extraction and validation hierarchy are stable.

The Domain shall also catalogue defence geometry and local exposure
evidence where available, but damage, material deformation, scour and
FSI remain stretch extensions after hydrodynamic load extraction is
stable. Source role: project-scope baseline and data-interface
decision. Supporting sources: event metadata from
\textcite{usgs2011tohokuEvent}, onshore observations from
\textcite{mori2012nationwide}, offshore observations from
\textcite{noaa2005dart} and \textcite{noaa2011japanDart}, and terrain
source candidates in \textcite{jodcJegg500}, \textcite{gebco2026grid},
\textcite{noaa2022etopo} and \textcite{gsiDemDownload}. Limitation:
the cited source records support data selection and handoff
requirements; they do not demonstrate that the project has executed
validation.

\section{Domain Deliverables}
\label{sec:res-dat-domain-deliverables}

% -------------------------------------------------------------------------
% WBS ID: RES-DAT-GEO
% Deliverable: Bathymetry, Topography and Coastal Geometry
% Status: In progress
% Evidence status: Regional model data requirements established; dataset selection pending
% Review status: Not reviewed
% Last updated: 2026-07-19
% Dependencies: RES-MOD-SPEC, RES-NUM-SPEC, RES-PHY-NRS
% Downstream interfaces: QGIS/geospatial domain lock, RES-VER-SPEC, Software data ingestion
% -------------------------------------------------------------------------

\section{RES-DAT-GEO --- Bathymetry, Topography and Coastal Geometry}
\label{sec:res-dat-geo}

\subsection{Purpose and Scope}
\label{sec:res-dat-geo-purpose}

This Deliverable defines the bathymetry, topography and coastal-geometry
data required to build the regional 2D NLSWE domain, the two
regional-to-local corridor extracts and the local 3D URANS--VOF
terrain surfaces. It covers dataset selection, fallback hierarchy,
coordinate reference systems, vertical datums, preprocessing and
handoff requirements. It does not define arbitrary corridor sketches:
the corridor footprint is a data/model interface constructed from the
event source, target coast, observation footprint, defence geometry and
numerical-boundary requirements.

\subsection{Research Questions}
\label{sec:res-dat-geo-research-questions}

The data decision asks:

\begin{enumerate}
  \item Which topobathymetric source is authoritative inside the Japan
    nearshore and local corridors?
  \item Which global products are acceptable for the outer regional
    domain or fallback coverage?
  \item How are horizontal CRS, vertical datum, sign convention,
    resolution and licence metadata preserved through clipping and
    interpolation?
  \item Which preprocessing uncertainty must be handed to `RES-VER`
    and Software before validation claims are made?
\end{enumerate}

\subsection{Terminology and Definitions}
\label{sec:res-dat-geo-definitions}

% TODO: Define the technical terminology, symbols, quantities and
% classifications used in this Deliverable.

\subsection{Literature Evidence}
\label{sec:res-dat-geo-literature-evidence}

% TODO: Record evidence from peer-reviewed papers, authoritative datasets,
% standards, technical reports and primary documentation.
%
% Do not create a detached literature-summary list. Explain how each source
% supports, contradicts or limits a project decision.

\textcite{jodcJegg500} documents J-EGG500 as a 500 m gridded
bathymetric dataset around Japan compiled from depth measurements from
Japanese marine survey institutions. Source role: supports the primary
Japan-region bathymetry candidate. Project-specific conclusion:
J-EGG500 is the preferred first candidate for offshore bathymetry in
the Kamaishi and Sendai corridors if licence, datum and import metadata
are acceptable. Limitation: the source page notes heterogeneous data
quality and smoothing, so local relief and harbour/breakwater details
shall not be assumed resolved.

\textcite{gebco2026grid} provides the current global 15 arc-second
terrain grid for ocean and land and a Type Identifier grid. Source
role: supports outer-domain and fallback topobathymetry. Project-
specific conclusion: GEBCO is selected as the first global fallback and
outer-domain source, not as the preferred source for local defence
geometry. Limitation: GEBCO's vertical datum is treated as mean sea
level while shallow-water source datums can vary, so nearshore datum
reconciliation remains mandatory.

\textcite{noaa2022etopo} provides NOAA ETOPO 2022 as a 15 arc-second
global relief model designed for applications including tsunami
forecasting and modelling. Source role: supports secondary global
fallback and cross-check. Project-specific conclusion: ETOPO shall be
used to fill or audit outer-domain coverage if J-EGG500/GEBCO coverage
or provenance is insufficient. Limitation: it is not selected for
local barrier-scale topography.

\textcite{gsiDemDownload} documents GSI digital elevation models with
1 m, 5 m and 10 m mesh products. Source role: supports Japan onshore
topography. Project-specific conclusion: GSI DEM is the preferred
onshore topographic source for corridor clipping and inundation/run-up
comparison where available. Limitation: mixed DEM classes have
different survey methods and accuracy; class boundaries shall be
retained in the provenance manifest.

\subsection{Governing Relations and Quantitative Definitions}
\label{sec:res-dat-geo-governing-relations}

% TODO: Record relevant equations, dimensionless groups, empirical models,
% extraction rules or quantitative definitions.
%
% Use align environments for displayed equations.
% Every displayed equation must have a unique \label{eq:...}.
% Do not place punctuation after displayed equations.

\subsection{Required Data Variables and Units}
\label{sec:res-dat-geo-required-data-variables-units}

% TODO: Complete this RES-DAT Domain-specific subsection.

\subsection{Candidate Data Sources}
\label{sec:res-dat-geo-candidate-data-sources}

The selected/fallback hierarchy for the Tōhoku case is:

\begin{enumerate}
  \item \textbf{Japan offshore bathymetry:} J-EGG500
    \parencite{jodcJegg500}, subject to licence, datum and coverage
    checks for the selected corridors.
  \item \textbf{Japan onshore topography:} GSI DEM
    \parencite{gsiDemDownload}, selecting the finest available mesh
    class that is legally usable and spatially continuous over each
    corridor.
  \item \textbf{Outer regional and fallback topobathymetry:}
    GEBCO\_2026 \parencite{gebco2026grid}.
  \item \textbf{Secondary global fallback or cross-check:}
    NOAA ETOPO 2022 \parencite{noaa2022etopo}.
\end{enumerate}

The selected approach is a stitched and provenance-preserving terrain
model rather than one unqualified raster. Each cell or clipped raster
tile shall retain the source product, source resolution, interpolation
method, horizontal CRS, vertical datum and sign convention.

\subsection{Spatial and Temporal Coverage}
\label{sec:res-dat-geo-spatial-temporal-coverage}

Coverage shall include:

\begin{itemize}
  \item the outer Tōhoku regional propagation domain, including
    artificial-boundary and sponge regions;
  \item the epicentre-to-Kamaishi corridor and enough lateral margin to
    contain the severe-impact observations, bay/coastline geometry and
    relaxation zones;
  \item the epicentre-to-Sendai coastal-plain corridor and enough
    inland coverage to include observed inundation limits;
  \item local 3D terrain patches around wall-type and
    obstacle/dissipating barrier classes.
\end{itemize}

Topobathymetric products are treated as static during one hydrodynamic
run. The earthquake source updates the bed through `RES-DAT-SRC`; it
does not change the static terrain-data selection.

\subsection{Coordinate and Datum Requirements}
\label{sec:res-dat-geo-coordinate-datum-requirements}

The regional model uses a horizontal projected coordinate frame
$(x,y)$ and a positive-upward bed elevation $b$. Source datasets shall
record horizontal CRS, vertical datum, units, coastline convention and
any tidal or geoid correction used before interpolation to the
finite-volume mesh.

Base acquisition coordinates shall remain longitude/latitude
$(\lambda,\phi)$ until the corridor target points are finalised in GIS.
The solver-facing terrain shall be delivered in a projected metric CRS
chosen for the final regional and local domains. Any dataset whose
horizontal CRS or vertical datum cannot be confirmed shall be marked
``usable only for visual screening'' until metadata are resolved.

\subsection{Data Processing and Transformation}
\label{sec:res-dat-geo-data-processing-transformation}

The required transformation pipeline is:

\begin{align}
  (\lambda,\phi,z)
  &\rightarrow
  (E,N,b_0)
  \rightarrow
  (x,y,b_0)
  \rightarrow
  b_0^{\mathrm{mesh}}
  \rightarrow
  b_{0,i}
  \rightarrow
  b_i(t)
  \label{eq:res-dat-geo-regional-bed-data-pipeline}
\end{align}

For positive-upward elevation data:

\begin{align}
  b_0
  &=
  z
  -
  z_{\mathrm{ref}}
  \label{eq:res-dat-geo-positive-upward-bed}
\end{align}

For raw depth $d$ positive downward:

\begin{align}
  b_0
  &=
  -d
  -
  z_{\mathrm{ref}}
  \label{eq:res-dat-geo-positive-downward-depth-conversion}
\end{align}

The cell-average static bed is:

\begin{align}
  b_{0,i}
  &=
  \frac{1}{A_i}
  \int_{K_i}
  b_0^{\mathrm{mesh}}(x,y)
  \,\dd A
  \label{eq:res-dat-geo-cell-average-static-bed}
\end{align}

QGIS/geospatial work owns dataset selection, domain placement,
coordinate transformation and source-data preparation. It is a
dependency of the regional model, not part of the mathematical model
itself.

\subsection{Uncertainty, Provenance and Licensing}
\label{sec:res-dat-geo-uncertainty-provenance-licensing}

The terrain manifest shall store, for every source tile or derived
terrain surface:

\begin{itemize}
  \item source product, provider, citation key and access URL;
  \item licence or usage note;
  \item horizontal CRS, vertical datum and sign convention;
  \item native resolution and interpolation or smoothing operations;
  \item clipping polygon, corridor identifier and mesh-projection
    method;
  \item uncertainty flags for datum conversion, coastline alignment,
    resampling, smoothing and source-age heterogeneity.
\end{itemize}

Project conclusion: a terrain raster without this metadata is not an
accepted validation input. Limitation: the metadata requirements do not
quantify the full bathymetric error; they make the uncertainty
auditable by `RES-VER`.

\subsection{Data Acceptance Criteria}
\label{sec:res-dat-geo-data-acceptance-criteria}

% TODO: Complete this RES-DAT Domain-specific subsection.

\subsection{Candidate Approaches}
\label{sec:res-dat-geo-candidate-approaches}

% TODO: Define the credible candidate approaches considered.

\subsection{Critical Comparison}
\label{sec:res-dat-geo-critical-comparison}

% TODO: Compare candidates using physically and project-relevant criteria.
% Do not select an approach solely on popularity or implementation ease.

\subsection{Assumptions and Applicability}
\label{sec:res-dat-geo-assumptions}

% TODO: State assumptions and the conditions under which they are valid.

\subsection{Limitations and Failure Conditions}
\label{sec:res-dat-geo-limitations}

% TODO: State where the evidence, model or selected approach ceases to be
% reliable or applicable.

\subsection{Project-Specific Conclusions}
\label{sec:res-dat-geo-conclusions}

The Tōhoku case-study terrain workflow shall use Japan-specific sources
where available and global products only where they are physically and
numerically appropriate. J-EGG500 and GSI DEM are selected as the
first candidates for Japan offshore and onshore data respectively.
GEBCO\_2026 is selected for outer-domain/fallback coverage, with ETOPO
2022 retained as a secondary fallback and consistency check. The local
barrier mesh shall not infer defence geometry from 15 arc-second global
products.

\subsection{Selected Approach and Rationale}
\label{sec:res-dat-geo-selected-approach}

Use a provenance-preserving terrain hierarchy:

\begin{quote}
  J-EGG500 offshore bathymetry $+$
  GSI onshore DEM where available, merged with GEBCO\_2026 for
  outer-domain and fallback coverage, and ETOPO 2022 as a secondary
  fallback/cross-check
\end{quote}

Source role: supports terrain-data finalisation for the regional model,
the two corridor extracts and local 3D terrain. Supporting citations:
\textcite{jodcJegg500}, \textcite{gsiDemDownload},
\textcite{gebco2026grid} and \textcite{noaa2022etopo}. Downstream
handoff: Software shall ingest these sources through a dataset
manifest, perform coordinate transformation, clip corridor rasters,
interpolate to mesh/field containers and write structured outputs with
source provenance.

\subsection{Unresolved Decisions}
\label{sec:res-dat-geo-unresolved-decisions}

% TODO: Record unresolved questions, required evidence and decision owners.

The mathematical data requirements are established. The actual QGIS
domain lock, final bathymetry/topography dataset, datum reconciliation,
coastline treatment, smoothing policy and mesh-projection workflow
remain open.

Open source-specific decisions are the final J-EGG500 and GSI licence
acceptance, any required Japanese geodetic/vertical datum conversion,
GEBCO/ETOPO blending boundary, and whether higher-resolution
site-specific harbour or defence surveys can be obtained for Kamaishi.
TODO--missing-source: no dedicated local Kamaishi bay-mouth breakwater
geometry or survey source has yet been confirmed in the WBS
bibliography.

\subsection{Requirements and Handoffs}
\label{sec:res-dat-geo-requirements}

% TODO: State requirements passed to other Research Domains, Software,
% verification activities or later execution work.

The data workstream shall provide mesh-projected static bed values,
source-bed increments where applicable, domain boundaries, sponge-layer
candidate regions and provenance metadata before validation or
implementation can be declared complete.

Software-facing requirements now begin in parallel with remaining
research: dataset manifest, coordinate-transform workflow,
corridor-polygon generation, raster/vector clipping,
bathymetry/topography interpolation, HDF5 or equivalent structured
terrain output, and replay-boundary metadata for the local 3D inlet.

\subsection{Deliverable Completion Record}
\label{sec:res-dat-geo-completion-record}

% TODO: Record whether the Deliverable is:
%
% - Not started
% - In progress
% - Draft complete
% - Reviewed
% - Accepted
%
% Acceptance requires:
%
% - the research questions to be answered;
% - claims to be supported by appropriate evidence;
% - assumptions and limitations to be explicit;
% - the project-specific conclusion to be stated;
% - downstream requirements to be traceable;
% - unresolved decisions to be closed or formally deferred.

% -------------------------------------------------------------------------
% WBS ID: RES-DAT-SRC
% Deliverable: Earthquake and Tsunami Source Data
% Status: In progress
% Evidence status: Regional model source-data requirements established; final dataset pending
% Review status: Not reviewed
% Last updated: 2026-07-19
% Dependencies: RES-PHY-SRC, RES-MOD-SRC, RES-MOD-SPEC
% Downstream interfaces: RES-NUM-SRC, RES-VER-SPEC, Software source ingestion
% -------------------------------------------------------------------------

\section{RES-DAT-SRC --- Earthquake and Tsunami Source Data}
\label{sec:res-dat-src}

\subsection{Purpose and Scope}
\label{sec:res-dat-src-purpose}

This Deliverable defines the earthquake event metadata, finite-fault
source candidates and bed-deformation data required to reconstruct the
2011 Tōhoku tsunami in the regional 2D NLSWE model. It covers source
data selection and projection into model-ready bed increments. It does
not choose a structural damage model, a local barrier geometry or a
solver implementation.

\subsection{Research Questions}
\label{sec:res-dat-src-research-questions}

The source-data decision asks:

\begin{enumerate}
  \item Which event metadata define the authoritative epicentre,
    origin time, magnitude and depth for corridor construction?
  \item Which finite-fault or tsunami-source reconstructions are
    candidates for regional bed forcing?
  \item Which observations were used to infer a source and therefore
    cannot later be claimed as independent validation evidence?
  \item Which data fields must Software receive to generate
    $\Delta b_m(x,y)$ and rupture-time functions reproducibly?
\end{enumerate}

\subsection{Terminology and Definitions}
\label{sec:res-dat-src-definitions}

% TODO: Define the technical terminology, symbols, quantities and
% classifications used in this Deliverable.

\subsection{Literature Evidence}
\label{sec:res-dat-src-literature-evidence}

% TODO: Record evidence from peer-reviewed papers, authoritative datasets,
% standards, technical reports and primary documentation.
%
% Do not create a detached literature-summary list. Explain how each source
% supports, contradicts or limits a project decision.

\textcite{usgs2011tohokuEvent} is retained as the event-metadata
source for the 2011 Great Tōhoku earthquake: origin time
2011-03-11 05:46:24 UTC, epicentral coordinates near
$38.297^{\circ}\mathrm{N}$, $142.373^{\circ}\mathrm{E}$, depth
$29.0~\mathrm{km}$ and moment magnitude $M_{\mathrm{w}}=9.1$. Source
role: supports event identity, epicentre coordinate and corridor
origin. Project conclusion: these metadata define
$P_{\mathrm{e}}=(\lambda_{\mathrm{e}},\phi_{\mathrm{e}})$ until a
GIS-approved event-source point is selected. Limitation: event
metadata alone do not define the spatial seabed displacement field.

\textcite{usgsFiniteFaults} and \textcite{hayes2011rapid} are retained
for the USGS finite-fault source-data family and rapid Tōhoku source
characterisation. Source role: supports candidate finite-fault model
fields: subfault geometry, slip, rake, rupture time and rise time.
Project conclusion: the USGS finite-fault family is a baseline source
candidate for regional moving-bed forcing. Limitation: the finite-fault
model must still be converted into hydrodynamic bed increments and
tested against observation sets not used in the inversion.

\textcite{fujii2011tsunamisource} and \textcite{wei2013japan} remain
retained source-reconstruction references. Source role: support and
limit the tsunami-source calibration choices used in Tōhoku hindcasts.
Project conclusion: these sources are comparison and provenance
anchors when deciding whether the USGS finite-fault source or an
alternative tsunami-inversion source better supports the selected
Kamaishi and Sendai corridors. Limitation: any observation record used
to tune a source shall be marked $\mathcal{D}_{\mathrm{cal}}$ in
`RES-DAT-OBS`.

\subsection{Governing Relations and Quantitative Definitions}
\label{sec:res-dat-src-governing-relations}

% TODO: Record relevant equations, dimensionless groups, empirical models,
% extraction rules or quantitative definitions.
%
% Use align environments for displayed equations.
% Every displayed equation must have a unique \label{eq:...}.
% Do not place punctuation after displayed equations.

\subsection{Required Data Variables and Units}
\label{sec:res-dat-src-required-data-variables-units}

The dynamic moving-bed regional source requires, for each source
component or subfault $m$:

\begin{itemize}
  \item geometry, location, depth, length and width;
  \item strike, dip and rake;
  \item slip or displacement amplitude;
  \item rupture timing and rise-time function $R_m(t)$;
  \item projected bed increment $\Delta b_m(x,y)$;
  \item provenance and uncertainty metadata.
\end{itemize}

\subsection{Candidate Data Sources}
\label{sec:res-dat-src-candidate-data-sources}

The current source candidates are:

\begin{itemize}
  \item USGS event metadata for origin time, magnitude, depth and
    epicentre \parencite{usgs2011tohokuEvent};
  \item USGS finite-fault products and rapid finite-fault source
    characterisation \parencite{usgsFiniteFaults,hayes2011rapid};
  \item Fujii--Satake source reconstruction
    \parencite{fujii2011tsunamisource};
  \item DART-constrained and nested-model source workflows described
    by \textcite{wei2013japan}.
\end{itemize}

The source-data table in `RES-DAT-CAS` records each candidate's role,
downstream deliverable and limitation.

\subsection{Spatial and Temporal Coverage}
\label{sec:res-dat-src-spatial-temporal-coverage}

Spatial coverage shall include the full rupture footprint and enough
surrounding seafloor for interpolation of dynamic bed increments onto
the regional mesh. Temporal coverage shall include the origin time,
subfault rupture onset and rise-time function for every source
component retained in the moving-bed source.

\subsection{Coordinate and Datum Requirements}
\label{sec:res-dat-src-coordinate-datum-requirements}

Raw source coordinates shall be retained in longitude/latitude and
source-native depth convention. The solver-facing product shall be
projected into the same horizontal metric frame as `RES-DAT-GEO` and
shall use positive-upward bed increments compatible with
\cref{eq:res-dat-src-stage-bed-update}. Any conversion from fault slip
to seabed displacement shall record the elastic model, coordinate
frame, vertical sign convention and interpolation method.

\subsection{Data Processing and Transformation}
\label{sec:res-dat-src-data-processing-transformation}

The hydrodynamic source data shall be delivered as bed increments
compatible with:

\begin{align}
  b_i(t_s)
  &=
  b_{0,i}
  +
  \sum_m
  \Delta b_{m,i}
  R_m(t_s)
  \label{eq:res-dat-src-stage-bed-update}
\end{align}

where $t_s$ is the relevant Runge--Kutta stage time. Static source
initialisation may be retained as a verification or fallback case, but
the regional model specification treats dynamic moving bathymetry as the
selected mathematical source representation.

\subsection{Uncertainty, Provenance and Licensing}
\label{sec:res-dat-src-uncertainty-provenance-licensing}

The source manifest shall identify whether each observation used to
create or tune a source is calibration, validation or diagnostic
evidence. Finite-fault uncertainty shall be retained at least as source
version, inversion method, data types used, subfault discretisation,
rupture timing assumptions and published update date where available.
Licence and access notes shall be recorded before any finite-fault data
are stored in a project dataset bundle.

\subsection{Data Acceptance Criteria}
\label{sec:res-dat-src-data-acceptance-criteria}

A source dataset is accepted for regional replay only when it supplies
or can defensibly derive:

\begin{itemize}
  \item event origin metadata and source version;
  \item subfault geometry, strike, dip, rake, slip and depth;
  \item rupture timing or a documented static-source fallback;
  \item projected bed increments $\Delta b_{m,i}$ on the regional
    mesh;
  \item calibration/validation-use classification for all observation
    records used in the source construction;
  \item provenance, licence and transformation metadata.
\end{itemize}

\subsection{Candidate Approaches}
\label{sec:res-dat-src-candidate-approaches}

% TODO: Define the credible candidate approaches considered.

\subsection{Critical Comparison}
\label{sec:res-dat-src-critical-comparison}

% TODO: Compare candidates using physically and project-relevant criteria.
% Do not select an approach solely on popularity or implementation ease.

\subsection{Assumptions and Applicability}
\label{sec:res-dat-src-assumptions}

% TODO: State assumptions and the conditions under which they are valid.

\subsection{Limitations and Failure Conditions}
\label{sec:res-dat-src-limitations}

% TODO: State where the evidence, model or selected approach ceases to be
% reliable or applicable.

\subsection{Project-Specific Conclusions}
\label{sec:res-dat-src-conclusions}

The Tōhoku source decision is source-data finalisation, not solver
implementation. The USGS event metadata define the current corridor
origin and event identity. USGS finite-fault, Hayes rapid source,
Fujii--Satake and DART-constrained source workflows remain source
candidates until their projected bed increments are compared against
the selected offshore, nearshore and terrestrial observation split.

\subsection{Selected Approach and Rationale}
\label{sec:res-dat-src-selected-approach}

Use a source-candidate manifest with USGS event metadata as the event
identity baseline and USGS finite-fault/Fujii/DART-constrained
reconstructions as candidate bed-forcing products. The selected
hydrodynamic source shall be the candidate that best supports the
Kamaishi and Sendai corridor validation split without reusing
calibration records as independent validation.

\subsection{Unresolved Decisions}
\label{sec:res-dat-src-unresolved-decisions}

% TODO: Record unresolved questions, required evidence and decision owners.

The source-data variables required by the mathematical model are
established. The exact final fault-source dataset, any comparison
against independent source reconstructions, final rupture-time
functions and the accepted source-to-bed projection remain open
data-selection tasks.

\subsection{Requirements and Handoffs}
\label{sec:res-dat-src-requirements}

% TODO: State requirements passed to other Research Domains, Software,
% verification activities or later execution work.

The selected source dataset shall be marked as calibration,
validation-support or diagnostic evidence so `RES-VER` can distinguish
source-fit records from independent validation records.

Software shall start a source-data ingestion path in parallel with the
remaining research: source manifest, fault-model parser, coordinate
projection, bed-increment rasterisation, structured source-output
schema and reproducible moving-bed replay fields for the regional
solver.

\subsection{Deliverable Completion Record}
\label{sec:res-dat-src-completion-record}

% TODO: Record whether the Deliverable is:
%
% - Not started
% - In progress
% - Draft complete
% - Reviewed
% - Accepted
%
% Acceptance requires:
%
% - the research questions to be answered;
% - claims to be supported by appropriate evidence;
% - assumptions and limitations to be explicit;
% - the project-specific conclusion to be stated;
% - downstream requirements to be traceable;
% - unresolved decisions to be closed or formally deferred.

% -------------------------------------------------------------------------
% WBS ID: RES-DAT-OBS
% Deliverable: Observational and Validation Data
% Status: In progress
% Evidence status: Tohoku observation hierarchy and principal data sources established
% Review status: Not reviewed
% Last updated: 2026-07-19
% Dependencies: Selected historical event, candidate-defence data feasibility
% Downstream interfaces: RES-DAT-SCN, RES-DAT-CAS, RES-ML-DAT
% -------------------------------------------------------------------------

\section{RES-DAT-OBS --- Observational and Validation Data}
\label{sec:res-dat-obs}

\subsection{Purpose and Scope}
\label{sec:res-dat-obs-purpose}

This Deliverable now covers both event-validation observations and
defence-specific evidence for documented structures exposed to the selected
event. In addition to tsunami water-level and inundation observations, the
catalogue shall include documented defence damage, survival, failure,
pre-event geometry, post-event condition, current or reconstructed geometry,
drawings, photographs, reconstruction records and local exposure evidence.

\subsection{Research Questions}
\label{sec:res-dat-obs-research-questions}

% TODO: State the questions that must be answered before the Deliverable
% can be accepted.

\subsection{Terminology and Definitions}
\label{sec:res-dat-obs-definitions}

% TODO: Define the technical terminology, symbols, quantities and
% classifications used in this Deliverable.

\subsection{Literature Evidence}
\label{sec:res-dat-obs-literature-evidence}

% TODO: Record evidence from peer-reviewed papers, authoritative datasets,
% standards, technical reports and primary documentation.
%
% Do not create a detached literature-summary list. Explain how each source
% supports, contradicts or limits a project decision.

\subsubsection{Mori et al. (2012)}
\label{sec:res-dat-obs-evidence-mori-2012}

\textcite{mori2012nationwide} documented the coordinated post-event survey of the 2011 Tohoku tsunami. The survey involved 299 researchers from 64 universities and institutes, covered approximately $2{,}000~\mathrm{km}$ of the Japanese coastline, and contained 5,214 surveyed locations by the end of December 2011. The principal measured quantities were inundation height, run-up height, and inundation distance.

The paper is retained as the principal nationwide observational reference for validating maximum onshore water levels and inundation extent. It is not sufficient for validating velocity, discharge, momentum flux, impact pressure, or time-dependent structural loading, as the field dataset primarily records maximum water levels and landward limits.

\subsubsection{Wei et al. (2013)}
\label{sec:res-dat-obs-evidence-wei-2013}

\textcite{wei2013japan} demonstrated an observation hierarchy in which
deep-ocean tsunami records constrained the source, while offshore GPS
buoys, nearshore gauges, field-survey heights, and mapped inundation
provided subsequent assessments of the regional and coastal model.

The paper is retained as evidence for separating source-calibration
records from independent offshore and terrestrial validation data.

\subsubsection{NOAA/NCEI DART and Tsunami Data}
\label{sec:res-dat-obs-evidence-noaa-ncei}

\textcite{noaa2005dart} documents NOAA/NCEI as the long-term archive
for DART ocean bottom pressure data, including netCDF and CSV products.
\textcite{noaa2011japanDart} provides event-specific DART records for
the March 11, 2011 Japan earthquake, including raw observations,
computed tides and residual values. Source role: supports offshore
time-series validation and source-calibration audit. Project-specific
conclusion: DART records are mandatory candidates for regional
propagation and arrival-time assessment, but any DART record used to
infer a source shall be marked calibration evidence.

\textcite{noaaNceiHistoricalTsunami} documents the NCEI/WDS Global
Historical Tsunami Database, including source-event and run-up records
with observed maximum water heights, arrival or travel times and
inundation information where available. Source role: supports event
inventory, run-up record discovery and cross-checking against the
Mori/JTS survey. Project-specific conclusion: the database is a
screening and provenance source, not a replacement for the
case-specific field-survey dataset selected for validation.

\subsection{Observation Hierarchy}
\label{sec:res-dat-obs-observation}

The benchmark observation set shall be organised as:

\begin{enumerate}
  \item geodetic and seismic measurements for earthquake-source
    definition;
  \item deep-ocean bottom-pressure records for source and regional
    propagation assessment;
  \item DART residual water-level time series for arrival time,
    amplitude, phase and source-calibration provenance;
  \item offshore GPS-wave-gauge records for nearshore transformation;
  \item coastal, harbour and tide gauges for local waveform,
    arrival-time and maximum-water-height assessment;
  \item satellite altimetry for spatial wave-front assessment;
  \item surveyed inundation depth and run-up;
  \item mapped inundation extent.
\end{enumerate}

Each dataset shall record:

\begin{itemize}
  \item spatial coordinates and reference datum;
  \item sampling interval;
  \item filtering and de-tiding procedure;
  \item measurement uncertainty;
  \item missing-data intervals;
  \item whether it was used for calibration, validation, or
    diagnosis;
  \item the model stage against which it is compared.
\end{itemize}

\subsection{Calibration, Validation and Diagnostic Classification}
\label{sec:res-dat-obs-data-classification}

Every observational record shall be classified as one of:

\begin{align}
  \mathcal{D}
  &=
  \begin{cases}
    \mathcal{D}_{\mathrm{cal}},
    &
    \text{used to infer or tune the source or model parameters}
    \\
    \mathcal{D}_{\mathrm{val}},
    &
    \text{withheld for independent assessment}
    \\
    \mathcal{D}_{\mathrm{diag}},
    &
    \text{used for qualitative or diagnostic comparison}
  \end{cases}
  \label{eq:res-dat-obs-data-classification}
\end{align}

Agreement with $\mathcal{D}_{\mathrm{cal}}$ shall be described as a
reconstruction fit rather than independent validation.

Where one observation record is divided temporally, the inversion
window and withheld waveform interval shall be identified separately.

\subsection{Governing Relations and Quantitative Definitions}
\label{sec:res-dat-obs-governing-relations}

% TODO: Record relevant equations, dimensionless groups, empirical models,
% extraction rules or quantitative definitions.
%
% Use align environments for displayed equations.
% Every displayed equation must have a unique \label{eq:...}.
% Do not place punctuation after displayed equations.

\subsection{Required Data Variables and Units}
\label{sec:res-dat-obs-required-data-variables-units}

% TODO: Complete this RES-DAT Domain-specific subsection.

The observation record shall distinguish the following quantities:
\begin{itemize}
  \item arrival time or travel time of the first observed wave;
  \item maximum water height at a gauge, buoy, survey or run-up
    location;
  \item tsunami height before or at shoreline arrival, in metres;

  \item inundation height at an onshore location, in metres above the selected
    sea-level datum;
  \item run-up height at the maximum landward extent, in metres
    above the selected sea-level datum;
  \item inundation distance from the shoreline,
    in kilometres or metres;
  \item observation latitude and longitude in decimal
    degrees;
  \item measurement method and instrument;
  \item surveyed trace type;

  \item tidal-correction method;
  \item survey date and, where available, survey
    time;
  \item observation-quality or reliability classification.
\end{itemize}
Tsunami height, inundation height, and run-up height shall not be treated as
interchangeable validation variables.

\subsection{Early and Late-Time Validation}
\label{sec:res-dat-obs-timevalidation}

The early coherent wave field shall be assessed deterministically using:

\begin{itemize}
  \item arrival time;
  \item amplitude;
  \item waveform phase;
  \item dominant period;
  \item wave-front position.
\end{itemize}

Fixed-point time-series records and spatial wave-front observations
shall be treated as complementary evidence rather than interchangeable
validation products.

The later reflected and scattered wave field may additionally require:

\begin{itemize}
  \item moving-window variance;
  \item spectral energy;
  \item waveform envelope;
  \item cumulative energy flux;
  \item cumulative force impulse;
  \item duration above a defined hazard threshold.
\end{itemize}

Statistical late-time agreement shall supplement rather than replace
deterministic diagnosis. A persistent phase or amplitude discrepancy
shall first be examined for errors in the source, bathymetry,
boundaries, or numerical formulation.

\subsection{Candidate Data Sources}
\label{sec:res-dat-obs-candidate-data-sources}

Candidate defence-evidence sources include official survey reports, coastal
engineering records, reconstruction drawings, site photographs, government or
municipal infrastructure records, post-disaster damage surveys, peer-reviewed
case studies and georeferenced imagery. Each source shall be classified by
authority, date, provenance, licence or access constraint, and whether it
supports geometry, exposure, damage, survival or reconstruction evidence.

Candidate event-observation sources are DART bottom-pressure records
\parencite{noaa2005dart,noaa2011japanDart}, tide/coastal/harbour-gauge
records identified through official archives, the NCEI/WDS Global
Historical Tsunami Database \parencite{noaaNceiHistoricalTsunami}, the
Mori/JTS field survey \parencite{mori2012nationwide}, and the
observation sets documented by \textcite{fine2013japan} and
\textcite{wei2013japan}. Tide-gauge and harbour-gauge records remain a
TODO--source-confirmation item until the specific station data and
licence metadata for the two corridors are selected.

\subsection{Spatial and Temporal Coverage}
\label{sec:res-dat-obs-spatial-temporal-coverage}

% TODO: Complete this RES-DAT Domain-specific subsection.

The survey extended approximately $2{,}000~\mathrm{km}$ along the Japanese coastline. Maximum run-up greater than $10~\mathrm{m}$ was reported along approximately $530~\mathrm{km}$ of coast, while run-up greater than $20~\mathrm{m}$ occurred along approximately $200~\mathrm{km}$. More than 1,000 observations were collected on the Sendai Plain, providing comparatively dense coverage for regional inundation assessment. A restricted zone of approximately $30~\mathrm{km}$ around the Fukushima Daiichi Nuclear Power Plant was not surveyed. Observation density also varied substantially between regions and bays.

\subsection{Coordinate and Datum Requirements}
\label{sec:res-dat-obs-coordinate-datum-requirements}

% TODO: Complete this RES-DAT Domain-specific subsection.

\subsection{Data Processing and Transformation}
\label{sec:res-dat-obs-data-processing-transformation}

% TODO: Complete this RES-DAT Domain-specific subsection.

Field measurements were corrected to remove the astronomical tide at the time of maximum tsunami elevation. Near-field corrections used the NAOTIDEJ astronomical tide database after many coastal gauges were destroyed. The time of maximum tsunami elevation was estimated using a nonlinear long-wave simulation based on the FS40v46 finite-fault source model. The resulting near-field observations are therefore not completely independent of the Fujii source reconstruction. This dependence shall be retained in the data-provenance record and considered when defining independent validation tests.

\subsection{Uncertainty, Provenance and Licensing}
\label{sec:res-dat-obs-uncertainty-provenance-licensing}

% TODO: Complete this RES-DAT Domain-specific subsection.

The comparison between the astronomical tide database and surviving gauges produced a mean absolute tidal-elevation difference of approximately $6.0~\mathrm{cm}$. The complete relative tidal correction had a reported mean absolute difference of approximately $17.3~\mathrm{cm}$ and a standard deviation of approximately $25.0~\mathrm{cm}$. These values represent only part of the total observation uncertainty. Additional uncertainty may arise from watermark interpretation, instrument choice, reference-datum conversion, survey timing, measurement density, and local spatial variability. A uniform pointwise uncertainty shall not be assigned without inspecting the metadata of the selected observations.

\subsection{Data Acceptance Criteria}
\label{sec:res-dat-obs-data-acceptance-criteria}

For G1 data feasibility, at least one defence shall have sufficiently
reconstructable location, orientation, pre-event geometry, current or
reconstructed geometry, material assumptions, foundation assumptions, observed
event exposure, and damage, survival or reconstruction evidence.

Observational data selected for benchmark validation shall be accepted only when:
\begin{itemize}
  \item the measured quantity is identified explicitly as tsunami height, inundation height, run-up height, or inundation distance;
  \item arrival time and maximum water height are retained when the
    source archive provides them;
  \item the horizontal coordinates and vertical datum are available;
  \item the tidal-correction procedure is documented;
  \item the observation is located within the resolved computational domain;
  \item the local topography and bathymetry are represented at a resolution consistent with the comparison;
  \item model output and observation refer to the same physical quantity;
  \item the observation is not used simultaneously for calibration and independent validation without explicit separation.
\end{itemize}

\subsection{Candidate Approaches}
\label{sec:res-dat-obs-candidate-approaches}

% TODO: Define the credible candidate approaches considered.

\subsection{Critical Comparison}
\label{sec:res-dat-obs-critical-comparison}

% TODO: Compare candidates using physically and project-relevant criteria.
% Do not select an approach solely on popularity or implementation ease.

\subsection{Assumptions and Applicability}
\label{sec:res-dat-obs-assumptions}

% TODO: State assumptions and the conditions under which they are valid.

\subsection{Limitations and Failure Conditions}
\label{sec:res-dat-obs-limitations}

% TODO: State where the evidence, model or selected approach ceases to be
% reliable or applicable.

\subsection{Data Selection Conclusions}
\label{sec:res-dat-obs-dataconclusion}

No single observation type is sufficient to validate the full model
chain.

Deep-ocean records constrain source directivity and offshore
propagation but do not establish accurate inundation. Inundation
outlines constrain the flooded region but may conceal errors in local
water level, velocity, and source amplitude. Satellite altimetry
provides spatial wave-front information but has a lower signal-to-noise
ratio than bottom-pressure records.

The benchmark shall therefore use complementary temporal, spatial,
offshore, nearshore, and terrestrial observations.

\subsection{Project-Specific Conclusions}
\label{sec:res-dat-obs-conclusions}

% TODO: Convert the evidence into conclusions specific to the Studentship
% project. Do not merely restate literature findings.

The Mori et al.\ dataset shall be used to validate maximum inundation height, run-up height, and inundation extent within a selected regional case study. The complete nationwide dataset shall not be used as one undifferentiated objective function, as observation density, coastal regime, and data quality vary substantially between locations. The Sendai Plain is retained as a candidate regional inundation benchmark owing to its dense measurement coverage. Sanriku and Kamaishi Bay are retained as candidate local-transformation and barrier-study regions, subject to the acquisition of higher-resolution bathymetric, topographic, and dynamic observational data.

For the two active corridors, observation selection shall produce
separate offshore, nearshore and terrestrial bundles. The Kamaishi
bundle shall emphasise Sanriku/ria-coast transformation, bay/defence
interaction and local barrier evidence. The Sendai coastal-plain bundle
shall emphasise dense inundation, run-up and inland-extent observations.
The bundles shall not be combined into one undifferentiated validation
score.

\subsection{Selected Approach and Rationale}
\label{sec:res-dat-obs-selected-approach}

% TODO: Record the selected approach only after sufficient evidence exists.
% State the alternatives rejected and the reasons for rejection.

\subsection{Unresolved Decisions}
\label{sec:res-dat-obs-unresolved-decisions}

% TODO: Record unresolved questions, required evidence and decision owners.

\subsection{Requirements and Handoffs}
\label{sec:res-dat-obs-requirements}

% TODO: State requirements passed to other Research Domains, Software,
% verification activities or later execution work.

The numerical-verification workstream shall define separate error metrics for maximum inundation height, run-up height, and inundation distance. The mathematical- and numerical-model workstreams shall produce outputs that correspond directly to these observed quantities. Separate dynamic datasets shall be obtained for validating velocity, discharge, momentum flux, pressure, and impact loading. The present survey dataset shall not be used as sole validation evidence for those quantities.

Software shall preserve observation records in a manifest that records
source archive, station or survey identifier, coordinates, vertical
datum, time base, sampling interval, de-tiding/filtering procedure,
calibration/validation/diagnostic classification, and linked model
output variable.

\subsection{Deliverable Completion Record}
\label{sec:res-dat-obs-completion-record}

% TODO: Record whether the Deliverable is:
%
% - Not started
% - In progress
% - Draft complete
% - Reviewed
% - Accepted
%
% Acceptance requires:
%
% - the research questions to be answered;
% - claims to be supported by appropriate evidence;
% - assumptions and limitations to be explicit;
% - the project-specific conclusion to be stated;
% - downstream requirements to be traceable;
% - unresolved decisions to be closed or formally deferred.

% -------------------------------------------------------------------------
% WBS ID: RES-DAT-SCN
% Deliverable: Scenario and Severity Definition
% Status: In progress
% Evidence status: Tohoku corridor scenario baseline established
% Review status: Not reviewed
% Last updated: 2026-07-19
% Dependencies: RES-DAT-OBS, RES-DAT-SRC, RES-DAT-GEO
% Downstream interfaces: RES-DAT-CAS, RES-ML-SUR
% -------------------------------------------------------------------------

\section{RES-DAT-SCN --- Scenario and Severity Definition}
\label{sec:res-dat-scn}

\subsection{Purpose and Scope}
\label{sec:res-dat-scn-purpose}

This Deliverable defines one primary historical tsunami event and controlled
physical variations around that event for impact testing. Unrelated tsunami
events shall not be added to the initial project scope.

The active scenario baseline is the 2011 Tōhoku event with two
regional-to-local trajectories: epicentre to the Kamaishi region and
epicentre to the Sendai coastal plain. Scenario definition owns the
evidence-based corridor footprint; it does not define arbitrary drawing
operations or solver implementation.

\subsection{Research Questions}
\label{sec:res-dat-scn-research-questions}

The scenario decision asks:

\begin{enumerate}
  \item How are the two active regional-to-local corridors defined from
    event and target coordinates?
  \item Which evidence determines corridor width, offshore length,
    inland length and sponge/relaxation margins?
  \item Which scenario changes are active impact-test variations and
    which are deferred geometry or ML-generation tasks?
\end{enumerate}

\subsection{Terminology and Definitions}
\label{sec:res-dat-scn-definitions}

% TODO: Define the technical terminology, symbols, quantities and
% classifications used in this Deliverable.

\subsection{Literature Evidence}
\label{sec:res-dat-scn-literature-evidence}

% TODO: Record evidence from peer-reviewed papers, authoritative datasets,
% standards, technical reports and primary documentation.
%
% Do not create a detached literature-summary list. Explain how each source
% supports, contradicts or limits a project decision.

\subsection{Governing Relations and Quantitative Definitions}
\label{sec:res-dat-scn-governing-relations}

% TODO: Record relevant equations, dimensionless groups, empirical models,
% extraction rules or quantitative definitions.
%
% Use align environments for displayed equations.
% Every displayed equation must have a unique \label{eq:...}.
% Do not place punctuation after displayed equations.

\subsubsection{Corridor Coordinate Basis}
\label{sec:res-dat-scn-corridor-coordinate-basis}

Base geographic coordinates are longitude/latitude. The event point is:

\begin{align}
  P_{\mathrm{e}}
  &=
  \left(
    \lambda_{\mathrm{e}},
    \phi_{\mathrm{e}}
  \right)
  \label{eq:res-dat-scn-event-point}
\end{align}

For a target coast point
$P_{\mathrm{t}}=(\lambda_{\mathrm{t}},\phi_{\mathrm{t}})$, a local
metric displacement vector shall be constructed after projection or
local tangent-plane conversion:

\begin{align}
  \Delta\mathbf{r}
  &=
  \left[
    \Delta E,\,
    \Delta N
  \right]^{\mathsf{T}}
  \label{eq:res-dat-scn-target-displacement}
\end{align}

The signed rotation from the latitude/north axis to the corridor
centreline is:

\begin{align}
  \theta
  &=
  \operatorname{atan2}
  \left(
    \Delta E,\,
    \Delta N
  \right)
  \label{eq:res-dat-scn-corridor-bearing-angle}
\end{align}

where positive angle is east of north under the selected projected
coordinate convention. The rotated local basis is:

\begin{align}
  \hat{\mathbf{e}}_{x'}
  &=
  \left[
    \sin\theta,\,
    \cos\theta
  \right]^{\mathsf{T}}
  \label{eq:res-dat-scn-corridor-centreline-basis}
  \\
  \hat{\mathbf{e}}_{y'}
  &=
  \left[
    \cos\theta,\,
    -\sin\theta
  \right]^{\mathsf{T}}
  \label{eq:res-dat-scn-corridor-width-basis}
\end{align}

For any projected point $\mathbf{r}$ measured from $P_{\mathrm{e}}$,
the corridor coordinates are:

\begin{align}
  x'
  &=
  \Delta\mathbf{r}
  \cdot
  \hat{\mathbf{e}}_{x'}
  \label{eq:res-dat-scn-corridor-x-prime}
  \\
  y'
  &=
  \Delta\mathbf{r}
  \cdot
  \hat{\mathbf{e}}_{y'}
  \label{eq:res-dat-scn-corridor-y-prime}
\end{align}

Here $x'$ is the corridor centreline/propagation direction and $y'$ is
the corridor-width direction. A constant-width baseline corridor is:

\begin{align}
  \mathcal{C}
  &=
  \left\{
    \mathbf{r}:
    -L_{\mathrm{pre}}
    \leq
    x'
    \leq
    L_{\mathrm{inland}},
    \quad
    \left|y'\right|
    \leq
    \frac{W}{2}
  \right\}
  \label{eq:res-dat-scn-constant-width-corridor}
\end{align}

where $W$ is corridor width, $L_{\mathrm{pre}}$ is offshore/upstream
pre-impact length and $L_{\mathrm{inland}}$ is inland/downstream
coverage. Sponge or relaxation margins are added outside
$\mathcal{C}$ according to `RES-DAT-GEO`, `RES-MOD-IBC` and
`RES-NUM-IBC`.

Decision trace: Source role: supports coordinate-transform and
corridor-polygon generation for the data/model interface. Supporting
citations: the event coordinate is anchored to
\textcite{usgs2011tohokuEvent}; corridor width and inland/offshore
extent shall be supported by topobathymetric sources
\parencite{jodcJegg500,gsiDemDownload,gebco2026grid} and severe-impact
observations \parencite{mori2012nationwide,noaaNceiHistoricalTsunami}.
Limitation: the angles below use current coordinate proxies and must be
recomputed from the final GIS-selected target coordinates.

\subsection{Required Data Variables and Units}
\label{sec:res-dat-scn-required-data-variables-units}

Required scenario variables are
$P_{\mathrm{e}}$, $P_{\mathrm{t}}$, $\theta$, $W$,
$L_{\mathrm{pre}}$, $L_{\mathrm{inland}}$, upstream/downstream/lateral
sponge widths, target-region identifier, observation-footprint polygon,
defence-geometry footprint and source-data version.

\subsection{Candidate Data Sources}
\label{sec:res-dat-scn-candidate-data-sources}

% TODO: Complete this RES-DAT Domain-specific subsection.

\subsection{Spatial and Temporal Coverage}
\label{sec:res-dat-scn-spatial-temporal-coverage}

The two active trajectories are:

\begin{description}
  \item[$T_{\mathrm{K}}$] epicentre $\rightarrow$ Kamaishi region,
    using the current coordinate proxy
    $\theta_{\mathrm{K}}\approx -21.1^{\circ}$.
  \item[$T_{\mathrm{S}}$] epicentre $\rightarrow$ Sendai coastal
    plain, using the current coordinate proxy
    $\theta_{\mathrm{S}}\approx -97.4^{\circ}$.
\end{description}

The angles are not final GIS bearings. They shall be recomputed from
the accepted target coordinates, selected projection and corridor
origin before clipping or mesh generation.

\subsection{Coordinate and Datum Requirements}
\label{sec:res-dat-scn-coordinate-datum-requirements}

% TODO: Complete this RES-DAT Domain-specific subsection.

\subsection{Data Processing and Transformation}
\label{sec:res-dat-scn-data-processing-transformation}

% TODO: Complete this RES-DAT Domain-specific subsection.

\subsection{Uncertainty, Provenance and Licensing}
\label{sec:res-dat-scn-uncertainty-provenance-licensing}

% TODO: Complete this RES-DAT Domain-specific subsection.

\subsection{Data Acceptance Criteria}
\label{sec:res-dat-scn-data-acceptance-criteria}

A corridor definition is accepted only when:

\begin{itemize}
  \item final $P_{\mathrm{e}}$ and $P_{\mathrm{t}}$ coordinates are
    recorded with source and CRS metadata;
  \item $\theta$ is recomputed from the final GIS coordinates;
  \item $W$, $L_{\mathrm{pre}}$, $L_{\mathrm{inland}}$ and sponge
    widths are justified by source footprint, severe-impact
    observations, defence geometry, bathymetry/topography coverage and
    numerical relaxation requirements;
  \item observation and terrain datasets fully cover the corridor plus
    buffer;
  \item constant-width corridor adequacy is documented for each target.
\end{itemize}

\subsection{Candidate Approaches}
\label{sec:res-dat-scn-candidate-approaches}

The active scenario approach is a validated event baseline plus bounded local
variations in wave height, inundation depth, velocity, direction, duration and
selected severity scaling. These variations support robustness and
reinforcement screening within the event-conditioned domain.

The baseline corridor is constant-width. A narrowing or bay-following
corridor is deferred unless GIS evidence shows that the narrowing
represents true physical bay, coastline or bathymetric constraints
rather than a drawing convenience.

\subsection{Critical Comparison}
\label{sec:res-dat-scn-critical-comparison}

% TODO: Compare candidates using physically and project-relevant criteria.
% Do not select an approach solely on popularity or implementation ease.

The constant-width corridor is selected for the baseline because it
creates a reproducible clipping and replay-boundary interface. It is
limited where a ria bay, harbour mouth, river channel or coastal plain
controls the flow in a way that a rectangular corridor would
misrepresent. In those cases the physical constraint shall be encoded
from GIS coastline/topography evidence, not manually narrowed for
visual neatness.

\subsection{Assumptions and Applicability}
\label{sec:res-dat-scn-assumptions}

% TODO: State assumptions and the conditions under which they are valid.

\subsection{Limitations and Failure Conditions}
\label{sec:res-dat-scn-limitations}

% TODO: State where the evidence, model or selected approach ceases to be
% reliable or applicable.

\subsection{Project-Specific Conclusions}
\label{sec:res-dat-scn-conclusions}

The initial project is a single-event reconstruction and controlled
impact-testing study, not a multi-event generalisation study.

The Tōhoku/local comparison is explicitly two-corridor: Kamaishi for a
Sanriku/ria and defence-interaction case, and Sendai coastal plain for
flat-plain inundation comparison. This pairing supports barrier-class
comparison within one event while avoiding claims of cross-event
generalisability.

\subsection{Selected Approach and Rationale}
\label{sec:res-dat-scn-selected-approach}

Use two evidence-defined, constant-width baseline corridors:
$T_{\mathrm{K}}$ and $T_{\mathrm{S}}$. Corridor parameters are data
products owned by RES-DAT, not solver constants. Source role: supports
regional-to-local data extraction and replay-boundary generation.
Downstream handoff: Software shall generate corridor polygons,
clip raster/vector data, interpolate terrain, extract observation
subsets and write replay-boundary definitions from the same manifest.

\subsection{Unresolved Decisions}
\label{sec:res-dat-scn-unresolved-decisions}

Open decisions are final GIS-selected target coordinates, final
projection, corridor widths, upstream/downstream lengths, sponge
widths, observation-footprint inclusion rules and whether a
physically-constrained non-rectangular local boundary is required near
Kamaishi Bay.

\subsection{Requirements and Handoffs}
\label{sec:res-dat-scn-requirements}

Software shall start the corridor workflow in parallel with remaining
research: coordinate-transform module, corridor-polygon generation,
manifest-driven clipping, local inlet section extraction and
replay-boundary metadata. `RES-VER-TST` shall use the same corridor
parameters in interface-location and one-way-versus-coupled tests.

\subsection{Deliverable Completion Record}
\label{sec:res-dat-scn-completion-record}

% TODO: Record whether the Deliverable is:
%
% - Not started
% - In progress
% - Draft complete
% - Reviewed
% - Accepted
%
% Acceptance requires:
%
% - the research questions to be answered;
% - claims to be supported by appropriate evidence;
% - assumptions and limitations to be explicit;
% - the project-specific conclusion to be stated;
% - downstream requirements to be traceable;
% - unresolved decisions to be closed or formally deferred.

% -------------------------------------------------------------------------
% WBS ID: RES-DAT-UNC
% Deliverable: Data Quality, Uncertainty and Provenance
% Status: In progress
% Evidence status: Data provenance and preprocessing uncertainty requirements established
% Review status: Not reviewed
% Last updated: 2026-07-19
% Dependencies: RES-DAT-GEO, RES-DAT-SRC, RES-DAT-OBS, RES-DAT-SCN
% Downstream interfaces: RES-DAT-CAS, RES-VER-MET, Software data ingestion
% -------------------------------------------------------------------------

\section{RES-DAT-UNC --- Data Quality, Uncertainty and Provenance}
\label{sec:res-dat-unc}

\subsection{Purpose and Scope}
\label{sec:res-dat-unc-purpose}

This Deliverable defines the provenance, CRS, datum, resolution,
licence and preprocessing uncertainty records required for all Tōhoku
case-study data. It covers uncertainty traceability for source,
topobathymetry, observations, corridor geometry and derived replay
boundaries. It does not quantify final validation tolerances; those
belong to `RES-VER-MET`.

\subsection{Research Questions}
\label{sec:res-dat-unc-research-questions}

The uncertainty decision asks:

\begin{enumerate}
  \item Which metadata are mandatory before a dataset can be ingested?
  \item Which preprocessing operations change the meaning or error
    structure of the data?
  \item Which uncertainty flags must travel into validation reports and
    local replay boundaries?
\end{enumerate}

\subsection{Terminology and Definitions}
\label{sec:res-dat-unc-definitions}

% TODO: Define the technical terminology, symbols, quantities and
% classifications used in this Deliverable.

\subsection{Literature Evidence}
\label{sec:res-dat-unc-literature-evidence}

% TODO: Record evidence from peer-reviewed papers, authoritative datasets,
% standards, technical reports and primary documentation.
%
% Do not create a detached literature-summary list. Explain how each source
% supports, contradicts or limits a project decision.

\subsection{Governing Relations and Quantitative Definitions}
\label{sec:res-dat-unc-governing-relations}

% TODO: Record relevant equations, dimensionless groups, empirical models,
% extraction rules or quantitative definitions.
%
% Use align environments for displayed equations.
% Every displayed equation must have a unique \label{eq:...}.
% Do not place punctuation after displayed equations.

Each data asset $d$ shall carry a provenance record:

\begin{align}
  \mathcal{P}_{d}
  &=
  \left\{
    s_d,\,
    c_d,\,
    \mathcal{R}_{h,d},\,
    \mathcal{R}_{v,d},\,
    \Delta_d,\,
    \ell_d,\,
    \mathcal{T}_d,\,
    \mathcal{U}_d
  \right\}
  \label{eq:res-dat-unc-provenance-record}
\end{align}

where $s_d$ is source/provider, $c_d$ is citation key,
$\mathcal{R}_{h,d}$ is horizontal CRS, $\mathcal{R}_{v,d}$ is vertical
datum, $\Delta_d$ is native resolution, $\ell_d$ is licence/usage note,
$\mathcal{T}_d$ is the ordered preprocessing transform and
$\mathcal{U}_d$ is the uncertainty flag set.

\subsection{Required Data Variables and Units}
\label{sec:res-dat-unc-required-data-variables-units}

Mandatory fields are dataset/source name, provider, citation key, URL,
spatial coverage, temporal coverage, native resolution, horizontal CRS,
vertical datum, licence/usage note, project role, downstream
deliverable, preprocessing required and limitation.

\subsection{Candidate Data Sources}
\label{sec:res-dat-unc-candidate-data-sources}

% TODO: Complete this RES-DAT Domain-specific subsection.

\subsection{Spatial and Temporal Coverage}
\label{sec:res-dat-unc-spatial-temporal-coverage}

% TODO: Complete this RES-DAT Domain-specific subsection.

\subsection{Coordinate and Datum Requirements}
\label{sec:res-dat-unc-coordinate-datum-requirements}

No dataset shall be projected, merged or compared until its coordinate
and datum state is known or explicitly marked unresolved. Where a
source lacks confirmed CRS or datum metadata, the derived product shall
be flagged:

\begin{align}
  a_d
  &=
  \mathrm{screening\ only}
  \label{eq:res-dat-unc-screening-only-flag}
\end{align}

until the metadata are resolved.

\subsection{Data Processing and Transformation}
\label{sec:res-dat-unc-data-processing-transformation}

Preprocessing operations requiring explicit audit are reprojection,
vertical datum conversion, sign-convention conversion, tidal
correction, de-tiding, filtering, smoothing, gap filling, clipping,
raster/vector overlay, interpolation, mesh projection, source-to-bed
projection and replay-boundary sampling.

\subsection{Uncertainty, Provenance and Licensing}
\label{sec:res-dat-unc-uncertainty-provenance-licensing}

The uncertainty flag set shall include, where applicable:

\begin{itemize}
  \item source-age or survey-date heterogeneity;
  \item native-resolution mismatch across merged products;
  \item unknown or mixed vertical datums;
  \item horizontal CRS conversion uncertainty;
  \item smoothing or interpolation of sharp bathymetric/defence
    features;
  \item tide correction or harmonic-analysis dependence;
  \item source-inversion dependence for observation products;
  \item licence restrictions preventing redistribution;
  \item missing station, survey or defence metadata.
\end{itemize}

Source role: supports validation traceability and software ingestion.
Supporting dataset citations are recorded in `RES-DAT-CAS`.
Limitation: these flags identify uncertainty sources; they do not
replace quantitative uncertainty propagation.

\subsection{Data Acceptance Criteria}
\label{sec:res-dat-unc-data-acceptance-criteria}

A data asset is accepted for validation or solver input only when its
provenance record satisfies:

\begin{align}
  \mathcal{P}_{d}
  &\supseteq
  \left\{
    s_d,\,
    c_d,\,
    \mathcal{R}_{h,d},\,
    \mathcal{R}_{v,d},\,
    \Delta_d,\,
    \ell_d,\,
    \mathcal{T}_d
  \right\}
  \label{eq:res-dat-unc-minimum-accepted-provenance}
\end{align}

Records missing any of these fields may be used only for exploratory
screening unless the missing field is formally waived with a documented
project risk.

\subsection{Candidate Approaches}
\label{sec:res-dat-unc-candidate-approaches}

% TODO: Define the credible candidate approaches considered.

\subsection{Critical Comparison}
\label{sec:res-dat-unc-critical-comparison}

% TODO: Compare candidates using physically and project-relevant criteria.
% Do not select an approach solely on popularity or implementation ease.

\subsection{Assumptions and Applicability}
\label{sec:res-dat-unc-assumptions}

% TODO: State assumptions and the conditions under which they are valid.

\subsection{Limitations and Failure Conditions}
\label{sec:res-dat-unc-limitations}

% TODO: State where the evidence, model or selected approach ceases to be
% reliable or applicable.

\subsection{Project-Specific Conclusions}
\label{sec:res-dat-unc-conclusions}

The Tōhoku corridor comparison is sensitive to datum, resolution and
preprocessing differences. A single unqualified terrain raster, source
file or observation table is therefore not a valid project input. The
minimum accepted object is a data asset plus its provenance record
\(\mathcal{P}_{d}\).

\subsection{Selected Approach and Rationale}
\label{sec:res-dat-unc-selected-approach}

Use a manifest-first data workflow. Every source and derived product is
registered before use, and every transformation adds a new derived
record rather than overwriting the source identity.

\subsection{Unresolved Decisions}
\label{sec:res-dat-unc-unresolved-decisions}

Open decisions are final licence acceptance for J-EGG500/GSI-derived
products, exact vertical datum transformations, quantitative terrain
uncertainty values, observation station filtering rules and the
threshold for treating a missing metadata field as blocking rather than
screening-only.

\subsection{Requirements and Handoffs}
\label{sec:res-dat-unc-requirements}

Software shall implement a dataset manifest, immutable source records,
derived-product records, coordinate-transform audit logs, HDF5 or
equivalent structured output metadata and validation-report provenance
exports. `RES-VER-MET` shall consume the same uncertainty flags when
deciding whether a comparison is admissible.

\subsection{Deliverable Completion Record}
\label{sec:res-dat-unc-completion-record}

% TODO: Record whether the Deliverable is:
%
% - Not started
% - In progress
% - Draft complete
% - Reviewed
% - Accepted
%
% Acceptance requires:
%
% - the research questions to be answered;
% - claims to be supported by appropriate evidence;
% - assumptions and limitations to be explicit;
% - the project-specific conclusion to be stated;
% - downstream requirements to be traceable;
% - unresolved decisions to be closed or formally deferred.

% -------------------------------------------------------------------------
% WBS ID: RES-DAT-CAS
% Deliverable: Integrated Tsunami Case-Study Specifications
% Status: In progress
% Evidence status: Integrated Tohoku data-source and corridor specification established
% Review status: Not reviewed
% Last updated: 2026-07-19
% Dependencies: RES-DAT-GEO, RES-DAT-SRC, RES-DAT-OBS, RES-DAT-SCN, RES-DAT-UNC, RES-GEO-SPEC
% Downstream interfaces: RES-VER-SPEC, Software data handling, local 3D replay boundary
% -------------------------------------------------------------------------

\section{RES-DAT-CAS --- Integrated Tsunami Case-Study Specifications}
\label{sec:res-dat-cas}

\subsection{Purpose and Scope}
\label{sec:res-dat-cas-purpose}

This Deliverable replaces the previous mandatory primary-plus-secondary-event
specification with one authoritative event-reconstruction case, one validated
event baseline, bounded impact-testing variants, and documented
defence-specific local cases.

The active integrated case is the 2011 Tōhoku event with two
regional-to-local corridors: Kamaishi and Sendai coastal plain. The
case supports impact modelling, local comparison and barrier-class
comparison. ML, surrogate optimisation and generative geometry are
deferred until the hydrodynamic data pipeline and validation hierarchy
are accepted.

\subsection{Research Questions}
\label{sec:res-dat-cas-research-questions}

The integrated case-study decision asks:

\begin{enumerate}
  \item Which data sources define the event, terrain, observations,
    corridors and optional sensitivities?
  \item Which source is selected, fallback, comparison-only,
    optional or missing?
  \item Which downstream deliverable and software asset consumes each
    source?
  \item Which limitations prevent a source from being used as
    validation or solver input?
\end{enumerate}

\subsection{Terminology and Definitions}
\label{sec:res-dat-cas-definitions}

% TODO: Define the technical terminology, symbols, quantities and
% classifications used in this Deliverable.

\subsection{Literature Evidence}
\label{sec:res-dat-cas-literature-evidence}

% TODO: Record evidence from peer-reviewed papers, authoritative datasets,
% standards, technical reports and primary documentation.
%
% Do not create a detached literature-summary list. Explain how each source
% supports, contradicts or limits a project decision.

The integrated case combines official data products and peer-reviewed
Tohoku evidence. Event identity is anchored to the USGS event record
\parencite{usgs2011tohokuEvent}. Candidate finite-fault and tsunami
source products are retained from USGS finite-fault documentation,
Hayes rapid source characterisation, Fujii--Satake and DART-constrained
hindcast evidence
\parencite{usgsFiniteFaults,hayes2011rapid,fujii2011tsunamisource,wei2013japan}.
Terrain sources are drawn from J-EGG500, GSI DEM, GEBCO\_2026 and ETOPO
2022 \parencite{jodcJegg500,gsiDemDownload,gebco2026grid,noaa2022etopo}.
Observation sources include NOAA/NCEI DART, the NCEI/WDS historical
tsunami database, Mori/JTS and published observation syntheses
\parencite{noaa2005dart,noaa2011japanDart,noaaNceiHistoricalTsunami,mori2012nationwide,fine2013japan}.
ERA5 and JRA-55 are retained only as optional wind/reanalysis
sensitivity sources, not baseline hydrodynamic forcing
\parencite{copernicusEra5,jmaJra55}.

\subsection{Governing Relations and Quantitative Definitions}
\label{sec:res-dat-cas-governing-relations}

% TODO: Record relevant equations, dimensionless groups, empirical models,
% extraction rules or quantitative definitions.
%
% Use align environments for displayed equations.
% Every displayed equation must have a unique \label{eq:...}.
% Do not place punctuation after displayed equations.

\subsection{Required Data Variables and Units}
\label{sec:res-dat-cas-required-data-variables-units}

The integrated case-study manifest shall include:

\begin{itemize}
  \item event metadata, finite-fault source version and derived bed
    increments;
  \item terrain source tiles and derived mesh-projected bed fields;
  \item observation bundles classified as calibration, validation or
    diagnostic;
  \item corridor polygons for $T_{\mathrm{K}}$ and $T_{\mathrm{S}}$;
  \item local barrier-class definitions and terrain/geometry patches;
  \item replay-boundary histories for local 3D inlet forcing;
  \item provenance records \(\mathcal{P}_d\) for every source and
    derived product.
\end{itemize}

\subsection{Candidate Data Sources}
\label{sec:res-dat-cas-candidate-data-sources}

\begin{landscape}
\scriptsize
\setlength{\tabcolsep}{3pt}
\begin{longtable}{
    p{0.09\linewidth}
    p{0.08\linewidth}
    p{0.07\linewidth}
    p{0.10\linewidth}
    p{0.17\linewidth}
    p{0.10\linewidth}
    p{0.14\linewidth}
    p{0.10\linewidth}
  }
  \caption{RES-DAT integrated data-source decision table for the Tōhoku corridor case}
  \label{tab:res-dat-cas-data-source-decision-table}
  \\
  \hline
  Dataset/source name
  &
  Provider
  &
  Citation key
  &
  URL
  &
  Coverage, resolution, CRS and vertical datum
  &
  Licence/usage note
  &
  Project role and downstream deliverable
  &
  Preprocessing required and limitation
  \\
  \hline
  \endfirsthead
  \hline
  Dataset/source name
  &
  Provider
  &
  Citation key
  &
  URL
  &
  Coverage, resolution, CRS and vertical datum
  &
  Licence/usage note
  &
  Project role and downstream deliverable
  &
  Preprocessing required and limitation
  \\
  \hline
  \endhead
  J-EGG500
  &
  Japan Oceanographic Data Center
  &
  \texttt{jodcJegg500}
  &
  \url{https://www.jodc.go.jp/jodcweb/JDOSS/infoJEGG.html}
  &
  Spatial: Japan surrounding seas. Temporal: static compiled bathymetry with heterogeneous survey dates. Resolution: 500 m grid. CRS/datum: confirm from delivered metadata before ingest; bathymetric sign convention requires conversion to positive-upward $b$.
  &
  Usage/licence to be confirmed from JODC terms before redistribution.
  &
  Preferred Japan offshore bathymetry candidate for `RES-DAT-GEO`, `RES-DAT-CAS`, regional mesh and corridor terrain.
  &
  Reproject, convert depth sign/datum, clip corridors, interpolate to meshes. Limitation: smoothing and mixed data quality may suppress harbour-scale relief.
  \\
  GEBCO\_2026 Grid
  &
  GEBCO Bathymetric Compilation Group
  &
  \texttt{gebco2026grid}
  &
  \url{https://www.gebco.net/data-products/gridded-bathymetry-data}
  &
  Spatial: global ocean/land terrain. Temporal: 2026 release. Resolution: 15 arc-second grid. CRS: geographic grid. Vertical datum: treated as mean sea level, with shallow-water datum caveat.
  &
  Public-domain use with required attribution/disclaimer.
  &
  Selected outer-domain and first fallback topobathymetry for `RES-DAT-GEO`, regional domain and source-to-coast propagation.
  &
  Clip outer domain, compare with J-EGG500/GSI boundaries, resample to mesh. Limitation: not sufficient for local defence geometry or high-resolution inundation features.
  \\
  NOAA ETOPO 2022
  &
  NOAA/NCEI
  &
  \texttt{noaa2022etopo}
  &
  \url{https://www.ncei.noaa.gov/products/etopo-global-relief-model}
  &
  Spatial: global relief. Temporal: 2022 release. Resolution: 15, 30 and 60 arc-second products. CRS/datum: confirm product metadata; elevation/bathymetry grid for global relief.
  &
  NOAA citation and DOI required; redistribution terms to be checked with source package.
  &
  Secondary global fallback and audit source for `RES-DAT-GEO` and outer-domain sensitivity.
  &
  Clip/resample and compare against GEBCO. Limitation: fallback only for local corridors and not baseline local terrain.
  \\
  GSI DEM
  &
  Geospatial Information Authority of Japan
  &
  \texttt{gsiDemDownload}
  &
  \url{https://service.gsi.go.jp/kiban/app/help/}
  &
  Spatial: Japan onshore. Temporal: source-age varies by DEM class. Resolution: 1 m, 5 m and 10 m mesh products. CRS/datum: confirm in JPGIS/GML metadata before ingest.
  &
  GSI usage terms and attribution required; redistribution status to be confirmed.
  &
  Preferred onshore topography for `RES-DAT-GEO`, run-up/inundation comparison and local terrain patches.
  &
  Convert JPGIS/GML, reproject, vertical-datum reconcile, clip corridors. Limitation: mixed DEM classes and coverage gaps require provenance flags.
  \\
  USGS Tōhoku event record
  &
  U.S. Geological Survey
  &
  \texttt{usgs2011tohokuEvent}
  &
  \url{https://earthquake.usgs.gov/earthquakes/eventpage/official20110311054624120_30/executive}
  &
  Spatial: epicentre near $38.297^{\circ}\mathrm{N}$, $142.373^{\circ}\mathrm{E}$. Temporal: 2011-03-11 05:46:24 UTC. Resolution: event metadata. CRS/datum: geographic coordinates, depth metadata.
  &
  USGS public information; cite source record.
  &
  Event identity and corridor origin for `RES-DAT-SRC` and `RES-DAT-SCN`.
  &
  Store immutable event metadata and convert to projected corridor origin. Limitation: does not define seabed displacement.
  \\
  USGS finite-fault family and Hayes rapid source
  &
  U.S. Geological Survey
  &
  \texttt{usgsFiniteFaults}; \texttt{hayes2011rapid}
  &
  \url{https://earthquake.usgs.gov/data/finitefault/}
  &
  Spatial: rupture/subfault footprint. Temporal: event rupture history. Resolution: subfault dependent. CRS/datum: source-model metadata required.
  &
  Cite USGS record and article; data-package usage terms to be checked.
  &
  Candidate source baseline for `RES-DAT-SRC` and moving-bed forcing.
  &
  Parse subfaults, project geometry, compute $\Delta b_m$, rasterise to regional mesh. Limitation: calibration observations must not be reused as independent validation.
  \\
  Fujii--Satake Tōhoku source
  &
  Fujii et al.
  &
  \texttt{fujii2011tsunamisource}
  &
  DOI in bibliography
  &
  Spatial: 2011 Tōhoku source region. Temporal: event source reconstruction. Resolution: source-model dependent. CRS/datum: paper/model metadata required.
  &
  Journal citation required.
  &
  Comparison source and provenance anchor for `RES-DAT-SRC` and observation de-tiding dependence.
  &
  Reconstruct source-to-bed projection if selected. Limitation: survey corrections in Mori et al. partly depend on FS40v46-style modelling.
  \\
  NOAA/NCEI DART ocean bottom pressure
  &
  NOAA/NCEI
  &
  \texttt{noaa2005dart}; \texttt{noaa2011japanDart}
  &
  \url{https://www.ncei.noaa.gov/products/natural-hazards/tsunamis-earthquakes-volcanoes/tsunamis/dart-ocean-bottom-pressure}
  &
  Spatial: DART station points, Pacific/global network. Temporal: event time series and archive records. Resolution: includes 15-second bottom-pressure records where available. CRS/datum: station metadata and pressure-to-water-level conversion required.
  &
  Cite NOAA/NCEI; data access unrestricted through archive tools unless station-specific terms state otherwise.
  &
  Offshore arrival, phase, amplitude and source-calibration audit for `RES-DAT-OBS` and `RES-VER-SPEC`.
  &
  Download station records, de-tide/filter, classify calibration/validation. Limitation: DART records used for source inference are calibration, not independent validation.
  \\
  NCEI/WDS Global Historical Tsunami Database
  &
  NOAA/NCEI/WDS
  &
  \texttt{noaaNceiHistoricalTsunami}
  &
  \url{https://www.ncei.noaa.gov/products/natural-hazards/tsunamis-earthquakes-volcanoes/tsunamis/global-historical-data}
  &
  Spatial: global tsunami source and run-up records. Temporal: historical-to-present database. Resolution: event/run-up record level. CRS/datum: per-record metadata; datum may be ambiguous.
  &
  DOI citation required; access unrestricted according to metadata.
  &
  Observation discovery, arrival-time and max-water-height screening for `RES-DAT-OBS`.
  &
  Query Tōhoku records, cross-check against field survey. Limitation: heterogeneous historical provenance; not sole validation source.
  \\
  Mori/JTS field survey
  &
  Mori, Takahashi and 2011 Tohoku Earthquake Tsunami Joint Survey Group
  &
  \texttt{mori2012nationwide}
  &
  DOI in bibliography
  &
  Spatial: approximately 2,000 km of Japanese coastline. Temporal: post-event 2011 survey. Resolution: surveyed point/extent records. CRS/datum: survey metadata and tide corrections required.
  &
  Journal citation required; raw-data access/licence to be confirmed.
  &
  Principal run-up, inundation-height and inundation-distance evidence for `RES-DAT-OBS`, Sendai and Sanriku/Kamaishi corridor validation.
  &
  Filter by corridor, datum-normalise, classify quality. Limitation: primarily maximum water-level and extent data, not velocity or pressure histories.
  \\
  Fine et al. observation/modelling synthesis
  &
  Fine et al.
  &
  \texttt{fine2013japan}
  &
  DOI in bibliography
  &
  Spatial: Japan 2011 tsunami observation/model context. Temporal: 2011 event. Resolution: paper-level synthesis. CRS/datum: underlying records must be obtained separately.
  &
  Journal citation required.
  &
  Supports observation hierarchy and regional propagation interpretation for `RES-DAT-OBS`.
  &
  Use as evidence, not raw data. Limitation: cannot replace station/survey archives.
  \\
  ERA5 and JRA-55
  &
  Copernicus Climate Change Service; Japan Meteorological Agency
  &
  \texttt{copernicusEra5}; \texttt{jmaJra55}
  &
  \url{https://climate.copernicus.eu/climate-reanalysis}; \url{https://www.data.jma.go.jp/jra/html/JRA-55/index_en.html}
  &
  Spatial: global atmospheric reanalysis. Temporal: multi-decadal products. Resolution: product dependent. CRS/datum: atmospheric grids.
  &
  Product-specific acknowledgement and terms required.
  &
  Optional wind/reanalysis sensitivity only; not baseline hydrodynamic forcing.
  &
  No baseline preprocessing. Limitation: not used unless a sensitivity question is opened and justified.
  \\
  Kamaishi bay-mouth breakwater defence records
  &
  TODO--provider
  &
  TODO--missing-source
  &
  TODO
  &
  Spatial: Kamaishi Bay and defence footprint. Temporal: pre-2011, 2011 damage/response and reconstruction states. Resolution/CRS/datum: unresolved.
  &
  Licence unresolved.
  &
  Required for real defence recreation and wall-type reference geometry in `RES-GEO` and local `RES-MOD`/`RES-VER`.
  &
  Locate authoritative drawings, surveys, damage reports or peer-reviewed defence papers. Limitation: no citation is created until source metadata are confirmed.
  \\
  \hline
\end{longtable}
\end{landscape}

\subsection{Spatial and Temporal Coverage}
\label{sec:res-dat-cas-spatial-temporal-coverage}

The integrated case covers the 2011 Tōhoku source-to-coast event and
two active regional-to-local corridor bundles:

\begin{enumerate}
  \item $T_{\mathrm{K}}$: epicentre to Kamaishi region, with
    Sanriku/ria, bay and defence-interaction emphasis;
  \item $T_{\mathrm{S}}$: epicentre to Sendai coastal plain, with
    coastal-plain inundation emphasis.
\end{enumerate}

The local 3D replay interval shall include all waves that materially
affect run-up, overtopping, transmitted/reflected energy and load
metrics, not only the first crest.

\subsection{Coordinate and Datum Requirements}
\label{sec:res-dat-cas-coordinate-datum-requirements}

All source coordinates are acquired as longitude/latitude and then
transformed into the selected projected metric CRS for regional and
local modelling. Vertical data shall be reconciled into the
positive-upward bed convention used by `RES-MOD` and `RES-NUM`.
Observation heights shall preserve their source datum and tide
correction before any comparison to simulated $\eta$, run-up or
inundation variables.

\subsection{Data Processing and Transformation}
\label{sec:res-dat-cas-data-processing-transformation}

The integrated data workflow is:

\begin{enumerate}
  \item register source assets in the dataset manifest;
  \item choose final $P_{\mathrm{e}}$, Kamaishi target and Sendai
    target coordinates;
  \item recompute $\theta_{\mathrm{K}}$ and $\theta_{\mathrm{S}}$;
  \item generate corridor polygons and buffer/sponge polygons;
  \item clip and reconcile terrain, source and observation datasets;
  \item interpolate terrain and bed increments to regional mesh
    fields;
  \item extract regional interface histories for local 3D replay;
  \item write structured HDF5 or equivalent outputs with provenance.
\end{enumerate}

\subsection{Uncertainty, Provenance and Licensing}
\label{sec:res-dat-cas-uncertainty-provenance-licensing}

The integrated case shall not be accepted unless every source in
\cref{tab:res-dat-cas-data-source-decision-table} has a provenance
record, licence/usage note and downstream role. Missing-source records
are allowed only as explicit blockers or TODOs; they shall not be
converted into fabricated citation keys.

\subsection{Data Acceptance Criteria}
\label{sec:res-dat-cas-data-acceptance-criteria}

An impact case is accepted only after event reconstruction validation and
defence-model credibility are established. The case record shall identify the
event baseline, local exposure, defence geometry version, material and
foundation assumptions, numerical model fidelity, convergence state, physical
validity state and uncertainty state.

For the current hydrodynamic baseline, an integrated case is accepted
for software implementation planning when it contains the dataset
manifest, coordinate transform, corridor polygons, clipped terrain,
source candidate records, observation bundles and local replay-boundary
format. This acceptance does not mean validation has been executed.

\subsection{Candidate Approaches}
\label{sec:res-dat-cas-candidate-approaches}

% TODO: Define the credible candidate approaches considered.

\subsection{Critical Comparison}
\label{sec:res-dat-cas-critical-comparison}

% TODO: Compare candidates using physically and project-relevant criteria.
% Do not select an approach solely on popularity or implementation ease.

\subsection{Assumptions and Applicability}
\label{sec:res-dat-cas-assumptions}

% TODO: State assumptions and the conditions under which they are valid.

\subsection{Limitations and Failure Conditions}
\label{sec:res-dat-cas-limitations}

% TODO: State where the evidence, model or selected approach ceases to be
% reliable or applicable.

\subsection{Project-Specific Conclusions}
\label{sec:res-dat-cas-conclusions}

The initial case-study specification is event-conditioned and
defence-specific. Secondary events and cross-event holdouts are deferred until
the single-event workflow is complete and reviewed.

The current integrated case is not a generic literature review. Each
source is retained because it supports, limits or defers a project
decision: terrain construction, source forcing, observation validation,
corridor definition, optional sensitivity or missing defence geometry.

\subsection{Selected Approach and Rationale}
\label{sec:res-dat-cas-selected-approach}

Use a manifest-driven Tōhoku case-study specification:

\begin{quote}
  USGS event identity and source candidates
  $+$
  J-EGG500/GSI/GEBCO/ETOPO terrain hierarchy
  $+$
  NOAA/NCEI and Mori/JTS observations
  $+$
  Kamaishi and Sendai corridor polygons
  $+$
  structured local replay-boundary outputs
\end{quote}

ERA5 and JRA-55 are explicitly optional sensitivity sources rather
than baseline forcing. Kamaishi defence records remain a missing-source
blocker for real defence recreation.

\subsection{Unresolved Decisions}
\label{sec:res-dat-cas-unresolved-decisions}

Open decisions are final GIS target coordinates, final terrain product
licence acceptance, final finite-fault/source selection, final
calibration/validation split, tide-gauge station selection, Kamaishi
defence-source acquisition, final corridor dimensions and final replay
boundary schema.

\subsection{Requirements and Handoffs}
\label{sec:res-dat-cas-requirements}

Data handling now begins in parallel with remaining research. Software
shall start with a modular MVP containing:

\begin{itemize}
  \item finite-volume data model and mesh/field containers;
  \item manifest-driven raster/vector I/O;
  \item coordinate-transform and corridor-polygon generation;
  \item bathymetry/topography interpolation and clipping;
  \item regional 2D scaffold and boundary framework;
  \item HDF5 or equivalent structured output schema;
  \item replay-boundary format for local 3D inlet forcing;
  \item local 3D scaffold that consumes replay histories without
    implementing the solver in this Research deliverable.
\end{itemize}

`RES-VER-SPEC` shall consume the same manifest when deciding which
records are calibration, validation or diagnostic evidence.

\subsection{Deliverable Completion Record}
\label{sec:res-dat-cas-completion-record}

% TODO: Record whether the Deliverable is:
%
% - Not started
% - In progress
% - Draft complete
% - Reviewed
% - Accepted
%
% Acceptance requires:
%
% - the research questions to be answered;
% - claims to be supported by appropriate evidence;
% - assumptions and limitations to be explicit;
% - the project-specific conclusion to be stated;
% - downstream requirements to be traceable;
% - unresolved decisions to be closed or formally deferred.


\clearpage
\printbibliography[
  heading=bibintoc,
  title={References}
]

\end{document}
